\begin{textsample}{POS Dim 1 – global\_south\_2025\_04 – Score 10.00 – global\_south\_2025\_04/123285230.txt}  \label{ex:f1_pos_329}
The uneasy ties between Palestinians and Arab states This is AI generated summarization, which may have errors.
For context, always refer to the full article.
LOOKING.
Palestinians react as they look at the site of an Israeli \textbf{strike} on a house, in Deir Al-Balah in the central Gaza Strip April 7, 2025.
REUTERS/Ramadan Abed Arab states express formal support for Palestinians, yet \textbf{regional} \textbf{tensions} and historical conflicts complicate their relationships with Palestinian communities across the Middle East DUBAI, UAE -- Arab states formally support the Palestinian cause but security, \textbf{regional} and sectarian \textbf{tensions} have strained their ties with the Palestinians scattered across the Middle East, leading to political and military crises.
Jordan Jordan was one of the countries to receive large numbers of Palestinian refugees when at least 700,000 Palestinians fled or were forced from their homes around Israel ' s foundation in 1948.
After Israel won the 1967 Arab-Israeli war, the Palestine Liberation Organization ( PLO ), headed by Yasser Arafat, moved to Jordan @ @ @ @ @ @ @ @ @ @ and threatened the rule of King Hussein.
King Hussein, who survived an \textbf{assassination} attempt after gunmen opened fire on his motorcade in 1970, hit back, and a civil war erupted.
Thousands of people were killed and the Palestinians were expelled to Beirut.
In 1994, Jordan became the second country after Egypt to sign a peace treaty that normalized ties with Israel.
Lebanon A multi-faceted civil war erupted in 1975 in a deeply sectarian Lebanon after the PLO resettled there from Jordan.
Palestinian refugee camps became a regular \textbf{target} of Israeli, \textbf{Syrian}, and \textbf{Lebanese} \textbf{militias} as they fought for control of the country once called the Switzerland of the Middle East.
In one of the worst atrocities of the 15-year conflict, \textbf{Lebanese} Christian militiamen backed by Israel massacred at least 800 Palestinian civilians at the Sabra and Shatila refugee camps.
While \textbf{Lebanese} authorities publicly advocate for the rights of Palestinian refugees, state policies that restrict their civil rights have been criticized by human rights organizations.
The @ @ @ @ @ @ @ @ @ @ permanent resettlement of Palestinians in Lebanon, thereby supporting their eventual return to their homeland.
Egypt The most populous Arab state has always regarded itself as a champion of the Palestinians.
Egypt has mediated between Israel and Hamas in the Gaza war that began in 2023, playing a similar role in past conflicts and peace negotiations.
Cairo is well connected to Palestinian factions including Hamas, which controlled Gaza from 2007.
However, Hamas is an offshoot of Egypt ' s Muslim Brotherhood, the Islamist group that President Abdel Fattah al-Sisi crushed after removing the Brotherhood ' s Mohamed Mursi from the presidency in 2013.
Egypt ' s priority has been to maintain security between Gaza and the Sinai Peninsula, where it has largely beaten back an insurgency.
Egypt helped Israel maintain a blockade of Gaza after Hamas took power.
Egypt was the first Arab state to make peace with Israel in 1979.
Palestinians living in Egypt say they faced increasing bureaucratic and security challenges from around this time.
United Arab Emirates The @ @ @ @ @ @ @ @ @ @ establish diplomatic ties with Israel in 30 years under the US-brokered Abraham Accords in 2020.
That broke a taboo on normalizing relations without the creation of a Palestinian state that had held since Jordan ' s peace deal.
Abu Dhabi has been a fierce opponent of Islamist groups across the \textbf{region}, including in Egypt, Sudan, and Libya.
Israel and the UAE have developed close economic and security ties, including defence cooperation.
Sudan Khartoum was long remembered by Israelis as the city where the Arab League in 1967 proclaimed its"Three No ' s"resolution on Israel -- no recognition, no peace and no negotiations.
But in October 2020 Sudan agreed to normalize relations with Israel, a decision driven by the prospects of economic relief, removal from the terrorism sponsors list, and broader international integration.
In return, US President Donald Trump ' s first administration agreed to remove Sudan from its list of state sponsors of terrorism, a designation that had isolated Sudan from the global economy. @ @ @ @ @ @ @ @ @ @ Israel.
Civilian groups, later sidelined by a coup and a civil war, were more reluctant.
Kuwait Relations between the PLO ' s Arafat and Kuwait were broken over his perceived sympathy for Iraqi President Saddam Hussein ' s decision to occupy his small Gulf Arab neighbour in 1990.
Kuwait had been an important supporter and funder of Arafat, who \textbf{launched} his Fatah movement while working there in 1964.
After the US-led war which defeated the Iraqis, hundreds of thousands of Palestinians were coerced into leaving or expelled outright, amid hostility from Kuwaitis who suspected them of disloyalty.
Iraq Under Saddam Hussein, Palestinians enjoyed subsidised housing, free education and the right to work -- rare privileges for foreigners.

% matched lemmas: assassination, launch, lebanese, militia, region, regional, strike, syrian, target, tension
\end{textsample}
